% ============================================================================
% CHAPTER 11: PRACTICAL NLP TASKS AND APPLICATIONS
% Bridges architecture knowledge to real-world NLP problem solving.
% ============================================================================

\chapter{Practical NLP Tasks and Applications}
\label{chap:practical_tasks}

\begin{tldr}
Knowing architectures is necessary but not sufficient---interviewers want to see that you can
\emph{apply} them to real tasks.  This chapter covers the major NLP application areas: text
classification, sentiment analysis, named entity recognition, relation extraction, question
answering, summarization, machine translation, and dialogue systems.  For each task, we cover
classical approaches, modern neural approaches, and LLM-based approaches, creating a
\textbf{decision framework} for when to use each.  The ability to select the right approach for a
given constraint set---data availability, latency budget, accuracy requirements, and
interpretability needs---is what separates a strong hire from a textbook reciter.
\end{tldr}

% ============================================================================
\section{Text Classification}
\label{sec:text_classification}
% ============================================================================

Text classification---assigning a label from a predefined set to a text input---is the most
common NLP task in production.  Spam filtering, content moderation, intent detection, topic
routing, and language identification are all text classification problems.  Despite its apparent
simplicity, the design space is rich and the ``right'' approach depends critically on
constraints.

\subsection{The Classical Pipeline}
\label{sec:classical_text_clf}

The classical text classification pipeline remains the right starting point for many production
systems:

\begin{enumerate}[label=\textbf{Step \arabic*.}]
    \item \textbf{Preprocessing:} Tokenize, lowercase, remove stopwords (optional),
          apply stemming or lemmatization (optional).
    \item \textbf{Feature extraction:} Convert text to a fixed-dimensional vector.
          TF-IDF over unigrams and bigrams is the workhorse.  Character n-grams
          help with misspellings and morphological variation.
    \item \textbf{Classification:} Train a supervised classifier on the feature vectors.
\end{enumerate}

\paragraph{Naive Bayes.}
Multinomial Naive Bayes models the likelihood of each word given the class using word
frequencies, then applies Bayes' theorem.  Despite the naive independence assumption
($P(x_1, \ldots, x_n \given y) = \prod_i P(x_i \given y)$), it works remarkably well for
text because:
\begin{itemize}
    \item High-dimensional sparse features are approximately conditionally independent given
          the class (the independence assumption is ``less wrong'' in high dimensions).
    \item The decision boundary depends only on the \emph{ratio} of class-conditional
          probabilities, so systematic errors in probability estimation often cancel.
    \item It requires very few parameters ($|V| \times |C|$), making it robust with small
          training sets.
\end{itemize}

\paragraph{Linear SVM.}
Support Vector Machines with a linear kernel are the classical SOTA for text classification.
Linear SVMs work exceptionally well on text because TF-IDF features are high-dimensional and
sparse---in this regime, data is often linearly separable, and the max-margin objective
generalizes well.  SVMs with TF-IDF bigrams remained competitive with early neural methods
on many benchmarks well into the deep learning era.

\paragraph{Logistic Regression.}
The probabilistic counterpart of linear SVMs.  Produces calibrated probabilities (useful when
you need a confidence score, not just a label), supports regularization (L1 for sparsity, L2
for generalization), and is trivially interpretable via feature weights.

\subsection{Deep Learning Approaches}
\label{sec:deep_text_clf}

\paragraph{TextCNN (Kim, 2014).}
Applies convolutional filters of varying widths (e.g., 3, 4, 5 tokens) over word embeddings,
followed by max-over-time pooling and a fully connected classifier.  Each filter width captures
n-gram patterns at a different scale.  TextCNN is fast, parallelizable, and surprisingly
effective---it was the first neural model to consistently beat SVMs on sentiment and topic
classification.

\paragraph{BiLSTM + Attention.}
A bidirectional LSTM encodes the sequence, and an attention mechanism computes a weighted sum
over hidden states to form a document representation.  The attention weights provide a form of
interpretability (which tokens mattered), though attention weights should not be over-interpreted
as explanations.

\paragraph{Fine-tuned BERT/RoBERTa.}
The modern default for text classification.  Add a linear classification head on the
\texttt{[CLS]} token and fine-tune the entire model (or use a frozen encoder with a trained
head for efficiency).  Advantages: captures deep contextual understanding, handles polysemy
and complex syntax, benefits from massive pretraining.  With as few as 1,000 labeled examples,
fine-tuned BERT typically outperforms all classical methods by a significant margin.

\paragraph{Few-Shot with LLMs.}
For prototyping or when labeled data is extremely scarce, prompting an LLM with a task
description and 2--8 examples can achieve competitive performance.  The tradeoff: much higher
latency and cost per prediction, no gradient-based learning from your data, and sensitivity
to prompt formatting.  Production systems typically use LLMs for labeling data to train a
smaller classifier, not for serving predictions directly.

\subsection{Multi-Label and Hierarchical Classification}

\paragraph{Multi-label classification.}
When each input can have multiple labels (e.g., a news article is both ``politics'' and
``economy''), the standard approaches are:
\begin{itemize}
    \item \textbf{Binary relevance:} Train an independent binary classifier per label.  Simple
          but ignores label correlations.
    \item \textbf{Classifier chains:} Train classifiers sequentially, with each model's input
          augmented by previous predictions.  Captures label dependencies at the cost of order
          sensitivity and error propagation.
    \item \textbf{Threshold tuning:} Use a single model with sigmoid outputs (not softmax) and
          tune the decision threshold per label on a validation set.  This is the practical
          default for BERT-based multi-label classifiers.
\end{itemize}

\paragraph{Hierarchical classification.}
When labels form a taxonomy (e.g., ``Sports $\to$ Football $\to$ Premier League''), flat
classification ignores the structure.  Taxonomy-aware approaches include: top-down classifiers
(predict coarse category first, then refine), hierarchical loss functions (penalize less for
predicting a sibling category than a distant one), and label embedding methods that encode the
taxonomy structure.

\subsection{Decision Framework: When Simple Models Win}

\begin{table}[H]
\centering
\caption{Text classification model selection guide.}
\label{tab:text_clf_selection}
\begin{tabular}{@{}lccccc@{}}
\toprule
\textbf{Scenario} & \textbf{NB/SVM} & \textbf{TextCNN} & \textbf{BERT} & \textbf{LLM} \\
\midrule
$<$ 100 labeled examples       & Best     & Poor   & Poor   & Good   \\
100--1K labeled examples       & Good     & Good   & Good   & Good   \\
$>$ 1K labeled examples        & Good     & Good   & Best   & Overkill \\
Latency $<$ 1 ms               & Best     & Good   & Poor   & Poor   \\
Latency $<$ 10 ms              & Best     & Best   & Good   & Poor   \\
Interpretability required      & Best     & Fair   & Poor   & Poor   \\
No GPU available               & Best     & Fair   & Poor   & Poor   \\
Rapidly changing label space   & Poor     & Poor   & Poor   & Best   \\
\bottomrule
\end{tabular}
\end{table}

\begin{keyinsight}[Always Start Simple]
The single most common mistake in NLP system design interviews is jumping straight to the most
complex model.  The correct answer almost always begins with: ``I would start with TF-IDF +
logistic regression as a baseline, establish performance, then only move to BERT or LLMs if the
baseline is insufficient and the latency/cost budget allows it.''  This signals engineering
maturity.  A Naive Bayes classifier trained in seconds can be serving traffic in an hour;
a fine-tuned BERT model requires GPU infrastructure, evaluation pipelines, and deployment
engineering.  The business value of shipping faster often outweighs the marginal accuracy gain.
\end{keyinsight}

% ============================================================================
\section{Sentiment Analysis}
\label{sec:sentiment_analysis}
% ============================================================================

Sentiment analysis determines the emotional polarity (positive, negative, neutral) of text.
While often framed as a simple classification task, production sentiment analysis requires
nuanced treatment of multiple levels of granularity.

\subsection{Levels of Sentiment Analysis}

\begin{enumerate}[label=\textbf{\arabic*.}]
    \item \textbf{Document-level:} Assign an overall sentiment to an entire review or article.
          This is standard text classification---TF-IDF + SVM or fine-tuned BERT work well.
    \item \textbf{Sentence-level:} Classify each sentence independently.  Useful when a
          document contains mixed sentiment (a positive review with one complaint).
    \item \textbf{Aspect-level:} Identify specific aspects (targets) and the sentiment toward
          each.  This is the most useful and most challenging level.
\end{enumerate}

\subsection{Aspect-Based Sentiment Analysis (ABSA)}

ABSA decomposes sentiment into structured triples: \textbf{(aspect, opinion, sentiment)}.
Consider the review:

\begin{quote}
\emph{``The food was excellent but the service was painfully slow.''}
\end{quote}

The desired output is:
\begin{itemize}
    \item (food, excellent, \textsc{positive})
    \item (service, painfully slow, \textsc{negative})
\end{itemize}

\paragraph{Subtasks of ABSA.}
\begin{itemize}
    \item \textbf{Aspect term extraction:} Identify the aspect phrases (``food,'' ``service'').
          This is a sequence labeling problem (BIO tagging).
    \item \textbf{Aspect category detection:} Map extracted aspects to predefined categories
          (e.g., ``food quality,'' ``service speed'').  This is classification.
    \item \textbf{Opinion term extraction:} Identify the opinion phrases (``excellent,''
          ``painfully slow'').  Also sequence labeling.
    \item \textbf{Sentiment classification:} Given an aspect-opinion pair, classify sentiment.
\end{itemize}

\paragraph{Modern approaches.}
Recent work frames ABSA as a generative task: given the review text, generate structured
triples directly using a seq2seq model (T5, BART) or an LLM with structured output.  This
avoids the error propagation of pipeline approaches and handles implicit aspects (where the
aspect is not explicitly mentioned but implied by the opinion).

\subsection{Challenges in Sentiment Analysis}

\begin{itemize}
    \item \textbf{Sarcasm:} ``Oh great, another meeting that could have been an email.''
          The words are positive but the sentiment is negative.  Sarcasm detection requires
          understanding pragmatic context, speaker intent, and often world knowledge.
    \item \textbf{Implicit sentiment:} ``The restaurant closed after two months'' conveys
          negative sentiment without any explicit opinion words.
    \item \textbf{Negation:} ``Not bad'' is positive; ``not particularly good'' is negative.
          Simple bag-of-words approaches fail on negation scope.
    \item \textbf{Domain shift:} A sentiment model trained on movie reviews performs poorly on
          product reviews because the vocabulary, style, and aspect space differ.
    \item \textbf{Comparative sentiment:} ``The camera is better than the iPhone but worse
          than the Pixel''---requires understanding comparison targets and relative polarity.
\end{itemize}

\begin{warningbox}
Sarcasm detection remains largely unsolved by NLP systems, including LLMs.  In interviews,
resist the urge to claim that ``a large enough model will handle sarcasm.''  The strongest
answer acknowledges that sarcasm often requires shared cultural context, knowledge of the
speaker, and situational awareness that is not present in text alone.  Multimodal cues
(tone of voice, facial expression) are often essential, which is why text-only sarcasm
detection has a natural accuracy ceiling.
\end{warningbox}

% ============================================================================
\section{Named Entity Recognition}
\label{sec:ner}
% ============================================================================

Named Entity Recognition (NER) identifies and classifies spans of text into predefined entity
types.  Standard entity types include Person (PER), Organization (ORG), Location (LOC), and
Miscellaneous (MISC) from the CoNLL-2003 shared task.  Domain-specific NER extends to medical
entities (diseases, drugs, dosages), financial entities (companies, monetary amounts, dates),
legal entities (laws, courts, parties), and many others.

\subsection{Tag Schemes}

NER is formulated as a sequence labeling problem where each token receives a tag.

\paragraph{BIO scheme.}
The most common scheme uses three tag types:
\begin{itemize}
    \item \textbf{B-\textit{type}}: Beginning of an entity of the given type.
    \item \textbf{I-\textit{type}}: Inside (continuation of) an entity.
    \item \textbf{O}: Outside any entity.
\end{itemize}

\paragraph{BIOES scheme.}
Extends BIO with two additional tags:
\begin{itemize}
    \item \textbf{E-\textit{type}}: End of a multi-token entity.
    \item \textbf{S-\textit{type}}: Single-token entity (both beginning and end).
\end{itemize}
BIOES is slightly more expressive than BIO and handles adjacent entities of the same type
more cleanly (e.g., ``[New York] [Times]'' as ORG-LOC vs.\ a single ORG).

\subsection{Approach Evolution}
\label{sec:ner_evolution}

\paragraph{Classical: CRF with hand-crafted features.}
The pre-neural standard used Conditional Random Fields with features including: word identity,
word shape (capitalization patterns, digit patterns), prefix/suffix n-grams, POS tags,
gazetteer membership (lists of known entity names), and window features (context words).  CRFs
model the \emph{joint} probability of the entire label sequence $P(\vy \given \vx)$, capturing
label-label transition dependencies that per-token classifiers miss.

\paragraph{Neural: BiLSTM-CRF (Lample et al., 2016).}
The breakthrough that eliminated manual feature engineering.  A bidirectional LSTM encodes
contextual token representations, and a CRF layer on top models label transitions.  The
BiLSTM captures input dependencies (context); the CRF captures output dependencies (valid
label sequences).

\paragraph{Modern: BERT + token classification.}
Fine-tune BERT with a linear layer (or CRF layer) on top of token representations.  Each
subword token receives a label, and predictions are aggregated to the word level (typically
using the first subword's prediction).  This is the current production standard for NER when
labeled data is available.

\paragraph{LLM-based: few-shot NER via prompting.}
Prompt an LLM with examples and ask it to extract entities.  Effective for prototyping and
low-resource scenarios, but unreliable for strict schema adherence and inconsistent on entity
boundary detection.  Best used to generate training data for a dedicated NER model.

\begin{keyinsight}[Why Add a CRF Layer?]
A per-token softmax classifier makes \emph{independent} predictions for each token---it can
produce invalid sequences like \texttt{O I-PER I-LOC} (an inside tag without a corresponding
beginning tag).  A CRF layer models transition probabilities between labels, learning that
\texttt{I-PER} can only follow \texttt{B-PER} or \texttt{I-PER}, never \texttt{B-LOC}.  The
Viterbi algorithm then finds the globally optimal label sequence in $O(T \times K^2)$ time,
where $T$ is the sequence length and $K$ is the number of labels.  Adding a CRF on top of
BERT typically improves NER F1 by 0.5--1.5 points, with the largest gains on entities with
strict structural constraints (nested types, multi-token entities).
\end{keyinsight}

\subsection{Entity Linking}
\label{sec:entity_linking}

NER alone identifies mentions (``Apple''); entity linking resolves them to knowledge base
entries (Apple~Inc.\ Q312 in Wikidata vs.\ the fruit).  The pipeline is:

\begin{enumerate}[label=\textbf{\arabic*.}]
    \item \textbf{Mention detection:} Identify entity spans (from NER or coreference).
    \item \textbf{Candidate generation:} Retrieve a candidate set of KB entities for each
          mention using alias tables, string matching, or dense retrieval.
    \item \textbf{Candidate ranking:} Score each candidate using context---a cross-encoder
          over (mention context, entity description) is the modern standard.
    \item \textbf{NIL detection:} Determine if the mention refers to an entity not in the KB
          (the ``NIL'' case).
\end{enumerate}

\subsection{NER Challenges}

\begin{itemize}
    \item \textbf{Nested entities:} ``[The [University of California] Berkeley]'' contains an
          ORG inside an ORG.  Flat sequence labeling cannot represent this.  Solutions include
          span-based models that enumerate and classify candidate spans, or layered/stacked
          NER.
    \item \textbf{Cross-domain transfer:} A NER model trained on news text performs poorly on
          biomedical text because entity types, naming conventions, and context patterns differ.
          Domain adaptation techniques include continued pretraining on domain text, few-shot
          fine-tuning, and gazetteer augmentation.
    \item \textbf{Noisy text:} Social media, OCR output, and speech transcriptions contain
          misspellings, non-standard capitalization, and abbreviations that break features
          dependent on surface form.
\end{itemize}

\subsection{NER Pipeline: From Text to Knowledge}

\begin{figure}[H]
\centering
\begin{tikzpicture}[
    node distance=0.6cm and 1.0cm,
    box/.style={rectangle, draw=deepblue, fill=lightblue, minimum width=2.4cm,
                minimum height=0.8cm, font=\small\sffamily, rounded corners=2pt,
                thick, align=center},
    databox/.style={rectangle, draw=gray!60, fill=white, minimum width=2.4cm,
                    minimum height=0.6cm, font=\scriptsize\sffamily, rounded corners=2pt,
                    align=center},
    arr/.style={-{Stealth[length=5pt]}, thick, deepblue!80},
    lbl/.style={font=\scriptsize\sffamily, text=gray!70!black}
]

% Pipeline stages
\node[box] (text) {Raw Text};
\node[box, right=1.2cm of text] (ner) {NER Model};
\node[box, right=1.2cm of ner] (link) {Entity Linker};
\node[box, right=1.2cm of link] (kg) {Knowledge\\Graph};

% Data flow
\node[databox, below=0.4cm of text] (d1) {``Tim Cook led\\Apple's keynote''};
\node[databox, below=0.4cm of ner] (d2) {[Tim Cook]\textsubscript{PER}\\{}[Apple]\textsubscript{ORG}};
\node[databox, below=0.4cm of link] (d3) {Tim Cook $\to$ Q484876\\Apple $\to$ Q312};
\node[databox, below=0.4cm of kg] (d4) {(Tim Cook, CEO\_of,\\Apple Inc.)};

% Arrows
\draw[arr] (text) -- (ner);
\draw[arr] (ner) -- (link);
\draw[arr] (link) -- (kg);

% Labels
\node[lbl, above=0.15cm of ner] {Sequence\\labeling};
\node[lbl, above=0.15cm of link] {Candidate\\ranking};
\node[lbl, above=0.15cm of kg] {Triple\\extraction};
\end{tikzpicture}
\caption{End-to-end NER pipeline from raw text to a populated knowledge graph.  Each stage
feeds into the next: NER identifies mentions, entity linking resolves them to KB entries, and
relation extraction populates the graph.}
\label{fig:ner_pipeline}
\end{figure}

% ============================================================================
\section{Relation Extraction}
\label{sec:relation_extraction}
% ============================================================================

Relation extraction identifies semantic relationships between entities in text.  Given a
sentence and two marked entities, the task is to classify the relationship (or determine that
no relation holds).  For example:

\begin{quote}
\emph{``\textbf{[Tim Cook]}\textsubscript{PER} is the CEO of
\textbf{[Apple]}\textsubscript{ORG}.''}
$\to$ (\textit{Tim Cook}, \textsc{ceo\_of}, \textit{Apple})
\end{quote}

\subsection{Supervised Relation Extraction}

\paragraph{Entity marker approach.}
The modern standard is to insert special tokens around entities in the input, then classify
using the \texttt{[CLS]} representation or the pooled entity representations.  For BERT:

\begin{lstlisting}[language=Python, caption={Entity marker input format for relation extraction.}]
# Input format with typed entity markers
input = "[CLS] <e1:PER> Tim Cook </e1:PER> is the CEO of "
        "<e2:ORG> Apple </e2:ORG> . [SEP]"
# Classify: [CLS] embedding -> linear -> relation labels
\end{lstlisting}

Variants include entity start markers (use the hidden state at the entity start token) and
typed markers (include the entity type in the marker).  Entity start markers typically
outperform \texttt{[CLS]}-based classification because they carry entity-specific information.

\subsection{Distant Supervision}
\label{sec:distant_supervision}

When labeled relation data is scarce, \textbf{distant supervision} (Mintz et al., 2009)
generates noisy training data by aligning a knowledge base with text:

\begin{quote}
\emph{If a knowledge base contains the triple (Tim Cook, CEO\_of, Apple), then every
sentence mentioning both ``Tim Cook'' and ``Apple'' is labeled as expressing the
\textsc{ceo\_of} relation.}
\end{quote}

The fundamental problem is that this assumption is wrong.  The sentence ``Tim Cook bought an
Apple product'' mentions both entities but does not express the CEO\_of relation.  This
creates \textbf{noisy labels} that degrade classifier performance.

\paragraph{Noise mitigation strategies.}
\begin{itemize}
    \item \textbf{Multi-instance learning:} Group all sentences mentioning the same entity
          pair into a ``bag'' and predict at the bag level, allowing the model to select
          informative instances and ignore noisy ones.
    \item \textbf{Denoising with attention:} Use attention over instances in a bag to
          downweight noisy sentences (Lin et al., 2016).
    \item \textbf{Human-in-the-loop:} Use distant supervision for initial labels, then
          refine with a small amount of human annotation on the noisiest examples.
\end{itemize}

\subsection{Open Information Extraction}

Unlike supervised RE, which classifies into a fixed relation schema, \textbf{Open IE} extracts
arbitrary (subject, predicate, object) triples from text without predefined relation types.
For example, from ``Einstein was born in Ulm in 1879,'' Open IE extracts:

\begin{itemize}
    \item (Einstein, was born in, Ulm)
    \item (Einstein, was born in, 1879)
\end{itemize}

Classical Open IE systems (ReVerb, OpenIE~5) use syntactic patterns.  Modern approaches use
LLMs with structured output to extract triples, achieving broader coverage but requiring
careful post-processing to normalize predicates and resolve entity references.

\subsection{Modern: LLM-Based Extraction}

LLMs can perform relation extraction through prompting with structured output requirements
(JSON, XML).  This approach excels when:
\begin{itemize}
    \item The relation schema is complex or evolving (no need to retrain).
    \item Few-shot examples are available but full training data is not.
    \item The extraction requires reasoning (implicit relations, multi-sentence evidence).
\end{itemize}

The tradeoff is cost and latency: an LLM call per sentence is orders of magnitude slower and
more expensive than a fine-tuned BERT classifier.  Production systems often use LLMs to
generate training data, then distill into a smaller model for serving.

% ============================================================================
\section{Question Answering}
\label{sec:question_answering}
% ============================================================================

Question answering (QA) systems return precise answers to natural language questions.  The
task taxonomy is rich, and the right architecture depends heavily on the QA type.

\subsection{Extractive QA}
\label{sec:extractive_qa}

In extractive QA, the answer is a contiguous span within a given context passage.  The model
predicts the start and end positions of the answer span.

\paragraph{BERT for extractive QA.}
The standard approach concatenates the question and context:

\[
\text{Input:} \quad \texttt{[CLS]}\; \underbrace{q_1 \ldots q_m}_{\text{question}}\;
\texttt{[SEP]}\; \underbrace{c_1 \ldots c_n}_{\text{context}}\; \texttt{[SEP]}
\]

Two linear layers project each token's hidden state to start and end logits:
\[
    P_{\text{start}}(i) = \softmax(\mW_s \vh_i + b_s), \qquad
    P_{\text{end}}(j) = \softmax(\mW_e \vh_j + b_e)
\]
The predicted answer span is $\argmax_{i \leq j} P_{\text{start}}(i) \times P_{\text{end}}(j)$,
typically with a constraint that $j - i < L_{\max}$ (maximum answer length).

\begin{mathresult}{Extractive QA Loss}
\[
    \loss = -\frac{1}{N}\sum_{k=1}^{N}\left[
        \log P_{\text{start}}(s^{(k)}) + \log P_{\text{end}}(e^{(k)})
    \right]
\]
where $s^{(k)}$ and $e^{(k)}$ are the ground-truth start and end positions for example $k$.
The loss decomposes into two independent cross-entropy terms, one for start position and one
for end position.
\end{mathresult}

\paragraph{Handling unanswerable questions.}
SQuAD~2.0 introduced questions that have no answer in the context.  The model must decide
whether to extract a span or output ``no answer.''  The standard approach treats the
\texttt{[CLS]} token as a ``null span'' candidate: if the \texttt{[CLS]} score exceeds all
span scores by a threshold $\tau$ (tuned on the validation set), the model predicts ``no
answer.''

\subsection{Extractive QA: Architecture}

\begin{figure}[H]
\centering
\begin{tikzpicture}[
    node distance=0.7cm,
    box/.style={rectangle, draw=deepblue, fill=lightblue, minimum width=1.6cm,
                minimum height=0.6cm, font=\scriptsize\sffamily, rounded corners=2pt,
                thick, align=center},
    tok/.style={rectangle, draw=gray!50, fill=white, minimum width=0.9cm,
                minimum height=0.5cm, font=\scriptsize\ttfamily, align=center},
    logit/.style={rectangle, draw=deepgreen!70, fill=lightgreen, minimum width=0.9cm,
                  minimum height=0.5cm, font=\scriptsize\sffamily, align=center},
    arr/.style={-{Stealth[length=4pt]}, thick, deepblue!70},
    brace/.style={decorate, decoration={brace, amplitude=5pt, mirror}, thick, gray!60}
]

% Token row
\node[tok] (cls) {\texttt{[CLS]}};
\node[tok, right=0.15cm of cls] (q1) {When};
\node[tok, right=0.15cm of q1] (q2) {was};
\node[tok, right=0.15cm of q2] (q3) {BERT};
\node[tok, right=0.15cm of q3] (sep1) {\texttt{[SEP]}};
\node[tok, right=0.15cm of sep1] (c1) {BERT};
\node[tok, right=0.15cm of c1] (c2) {was};
\node[tok, right=0.15cm of c2] (c3) {released};
\node[tok, right=0.15cm of c3] (c4) {in};
\node[tok, right=0.15cm of c4] (c5) {2018};
\node[tok, right=0.15cm of c5] (sep2) {\texttt{[SEP]}};

% BERT encoder
\node[box, above=0.8cm of q3, minimum width=10.8cm] (bert) {BERT Encoder (12/24 layers)};

% Hidden states
\node[logit, above=0.6cm of bert.west, anchor=west, xshift=0.5cm] (hcls) {$\vh_{\text{CLS}}$};
\node[logit, above=0.6cm of bert.east, anchor=east, xshift=-3.0cm] (h5) {$\vh_{\text{2018}}$};

% Arrows from tokens to BERT
\draw[arr] (cls.north) -- ++(0,0.35);
\draw[arr] (q3.north) -- ++(0,0.35);
\draw[arr] (c1.north) -- ++(0,0.35);
\draw[arr] (c5.north) -- ++(0,0.35);
\draw[arr] (sep2.north) -- ++(0,0.35);

% Linear layers
\node[box, above=0.7cm of hcls, fill=lightyellow, draw=orange!60] (null) {Null\\score};
\node[box, above=0.7cm of h5, fill=lightyellow, draw=orange!60] (span) {Start/End\\logits};

\draw[arr] (hcls) -- (null);
\draw[arr] (h5) -- (span);

% Braces
\draw[brace] (q1.south west) -- (q3.south east)
    node[midway, below=6pt, font=\scriptsize\sffamily, text=gray!70!black] {Question};
\draw[brace] (c1.south west) -- (c5.south east)
    node[midway, below=6pt, font=\scriptsize\sffamily, text=gray!70!black] {Context};
\end{tikzpicture}
\caption{BERT-based extractive QA.  The question and context are concatenated with separator
tokens.  Two linear layers predict start and end logits over context tokens.  The
\texttt{[CLS]} token's score serves as the ``no answer'' signal.}
\label{fig:extractive_qa}
\end{figure}

\subsection{Abstractive QA}

In abstractive QA, the model generates a free-form answer rather than extracting a span.
This is necessary when:
\begin{itemize}
    \item The answer requires synthesizing information from multiple parts of the context.
    \item The answer requires rephrasing or summarizing (``yes/no'' questions, aggregation).
    \item No single span contains the full answer.
\end{itemize}

Architectures include encoder-decoder models (T5, BART) and decoder-only LLMs.  Abstractive
QA is the natural formulation for RAG-based systems, where the generator must synthesize
information from multiple retrieved passages.

\subsection{Open-Domain QA}
\label{sec:open_domain_qa}

Open-domain QA answers questions without a pre-specified context.  The system must
\emph{retrieve} relevant passages from a large corpus, then \emph{read} them to extract or
generate an answer.

\paragraph{The retriever-reader pipeline.}
\begin{enumerate}[label=\textbf{\arabic*.}]
    \item \textbf{Retriever:} Given the question, retrieve the top-$k$ passages from a corpus.
          DPR (Dense Passage Retrieval) uses a dual-encoder BERT model trained on
          question-passage pairs with in-batch negatives.
    \item \textbf{Reader:} Given the question and retrieved passages, extract or generate the
          answer.  The reader can be extractive (select a span from the best passage) or
          abstractive (generate an answer conditioned on all passages).
\end{enumerate}

Modern open-domain QA systems are essentially RAG systems (Chapter~\ref{chap:practical_tasks}
refers to Chapter~9 for RAG details).  The key design decisions are: retrieval method (sparse
vs.\ dense vs.\ hybrid), number of passages to retrieve, how to present multiple passages to
the reader, and how to handle conflicting information across passages.

\subsection{Multi-Hop QA}

Multi-hop QA requires reasoning over multiple documents or passages to arrive at an answer.
For example:

\begin{quote}
\emph{``What is the capital of the country where the Eiffel Tower is located?''} \\
Hop 1: Eiffel Tower $\to$ located in France \\
Hop 2: France $\to$ capital is Paris
\end{quote}

Approaches include iterative retrieval (retrieve, read, generate a follow-up query, retrieve
again), graph-based reasoning (construct an entity graph across documents and reason over
paths), and chain-of-thought prompting with LLMs (which naturally decompose multi-hop
questions into reasoning steps).

\subsection{Table QA}

Answering questions from structured tables introduces unique challenges: understanding table
structure, aggregation operations (count, sum, average), and compositional queries.

\begin{itemize}
    \item \textbf{TAPAS (Google):} Extends BERT with additional embeddings for row/column
          position and trains on table-question pairs.  Predicts cell selection and aggregation
          operator.
    \item \textbf{Text-to-SQL:} Convert the natural language question to a SQL query, execute
          against a database.  Modern approaches use LLMs with schema-aware prompting.
          Spider and WikiSQL are standard benchmarks.
\end{itemize}

% ============================================================================
\section{Text Summarization}
\label{sec:summarization}
% ============================================================================

Text summarization condenses a source document into a shorter representation that preserves
the key information.  The two fundamental approaches---extractive and abstractive---have
different strengths, limitations, and failure modes.

\subsection{Extractive Summarization}

Extractive summarization selects a subset of sentences from the source document and presents
them as the summary.  No new text is generated.

\paragraph{TextRank.}
A graph-based unsupervised method inspired by PageRank.  Build a graph where sentences are
nodes and edge weights are sentence similarity (cosine similarity of TF-IDF or embedding
vectors).  Run PageRank to score sentence importance.  Select the top-$k$ sentences, ordered
by their position in the original document.  No training data required.

\paragraph{BertSumExt (Liu and Lapata, 2019).}
Fine-tune BERT for sentence-level binary classification: each sentence receives a score
indicating whether it should be included in the summary.  Inter-sentence attention (using
segment embeddings alternated per sentence) captures document-level context.

\paragraph{Advantages and limitations.}
Extractive methods are inherently faithful---every sentence in the summary is a verbatim copy
from the source, so hallucination is structurally impossible.  However, extractive summaries
are often choppy (sentences lack coherent flow), contain irrelevant details (an entire sentence
is selected even if only one clause is relevant), and cannot paraphrase or abstract.

\subsection{Abstractive Summarization}

Abstractive summarization generates new text that may paraphrase, compress, or fuse
information from the source.

\paragraph{Key models.}
\begin{itemize}
    \item \textbf{BART (Lewis et al., 2019):} Pretrained as a denoising autoencoder (corrupt
          text, reconstruct).  Strong at generation tasks including summarization.
    \item \textbf{Pegasus (Zhang et al., 2020):} Gap-sentence generation pretraining---mask
          entire sentences and predict them.  This pretraining objective directly mirrors
          extractive summarization, leading to excellent performance with few fine-tuning
          examples.
    \item \textbf{LLMs with prompting:} Provide the document and ask for a summary.  Highly
          flexible (can control length, style, focus), but expensive and prone to hallucination
          without careful grounding.
\end{itemize}

\subsection{Long Document Summarization}

When the source document exceeds the model's context window, strategies include:

\begin{itemize}
    \item \textbf{Truncation:} Simply truncate the input.  Surprisingly effective for news
          (the lead paragraph often contains the key information) but catastrophic for long
          reports or legal documents.
    \item \textbf{Chunking + map-reduce:} Summarize each chunk independently, then summarize
          the summaries.  Simple to implement but loses cross-chunk dependencies.
    \item \textbf{Hierarchical:} Encode chunks at a local level, then attend over chunk
          representations at a global level.  Architectures like Longformer and LED
          (Long-document Encoder-Decoder) support inputs of 4K--16K tokens.
    \item \textbf{Iterative refinement:} Generate an initial summary, then refine it by
          re-reading the source and correcting errors.
\end{itemize}

\subsection{The Faithfulness Problem}

Faithfulness---the degree to which the summary accurately represents the source---is the
central challenge of abstractive summarization.

\begin{itemize}
    \item \textbf{Intrinsic hallucination:} The summary contradicts information in the source.
          ``The company's revenue increased by 15\%'' when the source says 10\%.
    \item \textbf{Extrinsic hallucination:} The summary adds information not present in or
          inferable from the source.  ``The CEO expressed optimism'' when the source does not
          mention the CEO's sentiment.
\end{itemize}

\paragraph{Detection methods.}
\begin{itemize}
    \item \textbf{NLI-based:} Use a natural language inference model to check if the summary
          is entailed by the source.  If the NLI model predicts ``contradiction,'' the summary
          likely contains an intrinsic hallucination.
    \item \textbf{Question-based (QuestEval):} Generate questions from the summary, answer
          them using the source, and compare.  Mismatches indicate hallucination.
    \item \textbf{Factual consistency models:} Dedicated classifiers trained on hallucination
          annotations (FactCC, DAE).
\end{itemize}

\begin{keyinsight}[Faithfulness Over Fluency]
In production summarization, faithfulness is more important than fluency.  A choppy but
accurate extractive summary is preferable to a fluent but hallucinated abstractive one.
Interviewers test this because it reveals practical judgment: candidates who optimize for
ROUGE scores without discussing faithfulness are missing the metric that actually matters
for deployment.  ROUGE measures n-gram overlap, which correlates with fluency but is
\emph{completely blind to factual correctness}.
\end{keyinsight}

% ============================================================================
\section{Machine Translation}
\label{sec:machine_translation}
% ============================================================================

Machine translation (MT) converts text from a source language to a target language.  It is one
of the oldest NLP tasks (dating to the 1950s) and has been the driving application behind
several NLP breakthroughs, including the attention mechanism and the Transformer itself.

\subsection{Neural Machine Translation}

The standard NMT architecture is the Transformer encoder-decoder.  The encoder processes the
source sentence with bidirectional self-attention; the decoder generates the target sentence
autoregressively with causal self-attention and cross-attention to the encoder outputs.

\paragraph{Training.}
NMT models are trained on parallel corpora (aligned source-target sentence pairs) using
teacher forcing: at each decoding step, the ground-truth previous token is provided as input
(not the model's own prediction).  The loss is cross-entropy over the target vocabulary at
each position:

\[
    \loss = -\sum_{t=1}^{T} \log P(y_t \given y_{<t}, \vx)
\]

\paragraph{Label smoothing.}
Standard practice in NMT applies label smoothing ($\epsilon = 0.1$): instead of training
against hard one-hot targets, the target distribution places $1 - \epsilon$ on the correct
token and distributes $\epsilon$ uniformly across the vocabulary.  This reduces overconfidence
and improves BLEU by 0.5--1.0 points.

\subsection{Multilingual Translation}

\paragraph{Multilingual models.}
Rather than training separate models for each language pair, multilingual models handle
many languages in a single model:
\begin{itemize}
    \item \textbf{mBART (Liu et al., 2020):} Denoising autoencoder pretrained on monolingual
          data in 25 languages, then fine-tuned on parallel data.  Enables zero-shot
          translation between language pairs not seen during fine-tuning.
    \item \textbf{NLLB-200 (Costa-juss\`{a} et al., 2022):} ``No Language Left Behind.''
          Covers 200 languages with a single model, including many low-resource languages.
          Uses language-specific routing and data augmentation strategies.
\end{itemize}

\subsection{Low-Resource Translation Strategies}

When parallel data is scarce for a language pair, several strategies help:

\begin{itemize}
    \item \textbf{Back-translation (Sennrich et al., 2016):} Train a reverse model
          (target $\to$ source), use it to translate monolingual target data into synthetic
          source data, then add the synthetic pairs to training.  This is the single most
          effective data augmentation technique for NMT.
    \item \textbf{Pivot translation:} Translate through a high-resource intermediate language
          (e.g., translate Swahili $\to$ English $\to$ French instead of Swahili $\to$
          French directly).
    \item \textbf{Transfer from related languages:} Initialize a low-resource model from a
          high-resource related language (e.g., Portuguese $\to$ Galician).
    \item \textbf{Multilingual training:} Include the low-resource pair in a multilingual
          model that also trains on high-resource pairs.  Cross-lingual transfer improves
          low-resource performance.
\end{itemize}

\subsection{LLMs as Translators}

Large language models achieve competitive translation quality through few-shot prompting,
especially for high-resource language pairs.  They excel at:
\begin{itemize}
    \item Handling domain-specific or informal text (where NMT models trained on formal
          parallel data struggle).
    \item Providing explanations of translation choices (useful for language learning).
    \item Adapting to user preferences (formality level, regional variants) via instructions.
\end{itemize}

However, for high-volume production translation, dedicated NMT models remain preferred due to
their lower latency, lower cost per token, and more predictable quality.  LLMs also
underperform dedicated NMT for low-resource languages where the LLM's pretraining data
is limited.

\subsection{MT Evaluation}

Evaluation is covered in depth in Chapter~10.  The key metrics in brief:

\begin{itemize}
    \item \textbf{BLEU:} N-gram precision with brevity penalty.  The traditional standard but
          increasingly recognized as inadequate---it penalizes valid paraphrases and correlates
          poorly with human judgment for high-quality systems.
    \item \textbf{COMET:} A learned neural metric that uses reference translations and source
          text.  Correlates much better with human judgment than BLEU.
    \item \textbf{Human evaluation:} The gold standard.  MQM (Multidimensional Quality Metrics)
          provides a structured framework for annotating translation errors by type and severity.
\end{itemize}

% ============================================================================
\section{Dialogue Systems}
\label{sec:dialogue_systems}
% ============================================================================

Dialogue systems engage in multi-turn conversation with users.  The two primary categories
---task-oriented and open-domain---have fundamentally different architectures, evaluation
criteria, and deployment challenges.

\subsection{Task-Oriented Dialogue}
\label{sec:task_oriented_dialogue}

Task-oriented dialogue systems help users accomplish specific goals (booking a restaurant,
checking flight status, troubleshooting a device).  The classical architecture is a pipeline
of specialized components.

\begin{figure}[H]
\centering
\begin{tikzpicture}[
    node distance=0.4cm and 0.8cm,
    box/.style={rectangle, draw=deepblue, fill=lightblue, minimum width=2.6cm,
                minimum height=1.0cm, font=\small\sffamily, rounded corners=3pt,
                thick, align=center},
    arr/.style={-{Stealth[length=5pt]}, thick, deepblue!80},
    lbl/.style={font=\scriptsize\sffamily, text=gray!70!black, align=center},
    databox/.style={rectangle, draw=gray!50, fill=white, minimum width=2.6cm,
                    minimum height=0.5cm, font=\scriptsize\sffamily, rounded corners=2pt,
                    align=center}
]

% Pipeline boxes
\node[box] (nlu) {NLU\\{\scriptsize Intent +\\Slot Filling}};
\node[box, right=1.0cm of nlu] (dst) {Dialogue State\\Tracker};
\node[box, right=1.0cm of dst] (policy) {Dialogue\\Policy};
\node[box, right=1.0cm of policy] (nlg) {NLG\\{\scriptsize Response\\Generation}};

% User input/output
\node[databox, left=0.8cm of nlu] (user) {User utterance};
\node[databox, right=0.8cm of nlg] (resp) {System response};

% Arrows
\draw[arr] (user) -- (nlu);
\draw[arr] (nlu) -- (dst);
\draw[arr] (dst) -- (policy);
\draw[arr] (policy) -- (nlg);
\draw[arr] (nlg) -- (resp);

% Example data below
\node[lbl, below=0.3cm of nlu] {Intent: \texttt{book\_table}\\Slots: \{time: 7pm,\\party\_size: 4\}};
\node[lbl, below=0.3cm of dst] {Belief state:\\cuisine=Italian\\time=7pm\\party\_size=4};
\node[lbl, below=0.3cm of policy] {Action:\\\texttt{confirm}\\+ \texttt{ask\_location}};
\node[lbl, below=0.3cm of nlg] {``A table for 4 at\\7pm. Which area\\do you prefer?''};

% DB arrow
\node[box, below=1.8cm of policy, fill=lightyellow, draw=orange!60, minimum width=2.0cm] (db) {Backend/DB};
\draw[arr, orange!60] (policy) -- (db);
\draw[arr, orange!60] (db) -- (policy);
\end{tikzpicture}
\caption{Classical task-oriented dialogue pipeline.  Each component is a specialized module:
NLU parses user intent and slots, the dialogue state tracker maintains the belief state
across turns, the policy decides the next system action (potentially querying a backend), and
NLG generates the response.}
\label{fig:tod_pipeline}
\end{figure}

\paragraph{Natural Language Understanding (NLU).}
The NLU module performs two tasks jointly:
\begin{itemize}
    \item \textbf{Intent classification:} Determine the user's goal (e.g., \texttt{book\_table},
          \texttt{cancel\_reservation}, \texttt{check\_hours}).  Standard text classification.
    \item \textbf{Slot filling:} Extract relevant parameters as named slots (e.g.,
          \texttt{time=7pm}, \texttt{party\_size=4}).  Sequence labeling with BIO tags.
\end{itemize}
Joint intent-slot models (where the intent classifier and slot tagger share a BERT encoder)
outperform separate models because intent and slots are mutually informative.

\paragraph{Dialogue State Tracking (DST).}
The DST module maintains a structured representation of the user's goals across turns.  The
\textbf{belief state} is a set of slot-value pairs accumulated from conversation history.
Key challenge: the user may correct, update, or partially specify values across turns.

\begin{quote}
\emph{Turn 1: ``I'd like Italian food.''} $\to$ \{cuisine: Italian\} \\
\emph{Turn 2: ``Actually, make it Chinese.''} $\to$ \{cuisine: Chinese\} \\
\emph{Turn 3: ``For 4 people at 7pm.''} $\to$ \{cuisine: Chinese, party\_size: 4, time: 7pm\}
\end{quote}

Modern DST models use BERT-based classifiers over the full dialogue history or generative
models that produce the belief state as text.

\paragraph{Dialogue Policy.}
Decides what action the system should take given the current belief state: ask for a missing
slot, confirm a value, query the database, or provide information.  Classical approaches use
handcrafted rules; modern approaches use reinforcement learning or supervised learning from
human demonstrations.

\paragraph{Natural Language Generation (NLG).}
Converts the system action into natural language.  Template-based NLG (``A table for
\{\texttt{party\_size}\} at \{\texttt{time}\}'') is reliable and predictable.  Neural NLG
(fine-tuned LMs) is more fluent but risks generating incorrect slot values.

\subsection{Open-Domain Dialogue}

Open-domain (chitchat) systems aim to engage users in free-form conversation without a
specific task goal.  Quality is measured by engagement, coherence, personality consistency,
and safety.

\paragraph{Challenges.}
\begin{itemize}
    \item \textbf{Consistency:} The system should not contradict itself across turns (``I'm
          from New York'' in turn 3, ``I grew up in London'' in turn 7).
    \item \textbf{Engagement:} Avoiding generic responses (``That's interesting.'') while
          maintaining relevance.
    \item \textbf{Grounding:} Connecting responses to verifiable facts (reducing hallucination).
    \item \textbf{Safety:} Avoiding harmful, biased, or inappropriate responses.
\end{itemize}

\subsection{Modern LLM-Based Dialogue}

Modern dialogue systems are increasingly built on LLMs with the following components:

\begin{itemize}
    \item \textbf{System prompt:} Defines persona, constraints, capabilities, and safety rules.
    \item \textbf{Conversation memory:} The full conversation history (or a summarized version
          for long conversations) is included in the prompt.
    \item \textbf{Tool integration:} The LLM can call external tools (search, calculator,
          APIs) using function calling or ReAct-style prompting.
    \item \textbf{Safety guardrails:} Input classifiers (detect prompt injection, harmful
          requests) and output classifiers (detect harmful, off-topic, or policy-violating
          responses).
    \item \textbf{Retrieval augmentation:} Ground responses in retrieved documents or
          knowledge base entries to reduce hallucination.
\end{itemize}

This architecture collapses the classical pipeline into a single LLM that handles NLU, DST,
policy, and NLG jointly.  The tradeoff is reduced controllability: when the LLM makes an
error in any stage, it is harder to diagnose and fix than in a modular pipeline.

% ============================================================================
\section{Information Extraction Pipeline}
\label{sec:ie_pipeline}
% ============================================================================

Information extraction (IE) is the umbrella task of converting unstructured text into
structured representations.  In practice, IE systems chain multiple NLP components into an
end-to-end pipeline.

\subsection{From Text to Knowledge Graph}

The canonical IE pipeline proceeds in stages:

\begin{enumerate}[label=\textbf{\arabic*.}]
    \item \textbf{Entity extraction:} NER identifies entity mentions in text.
    \item \textbf{Coreference resolution:} Cluster mentions that refer to the same entity
          (``Tim Cook,'' ``he,'' ``the Apple CEO'' $\to$ same person).
    \item \textbf{Relation extraction:} For each pair of (co-referenced) entities, classify
          the relationship.
    \item \textbf{Entity linking:} Resolve each entity cluster to a canonical knowledge base
          entry.
    \item \textbf{Knowledge graph population:} Insert the extracted triples
          (subject, relation, object) into a knowledge graph.
\end{enumerate}

Each stage has error rates that compound: if NER has 90\% recall and RE has 85\% accuracy,
the end-to-end pipeline's effective accuracy is at most $0.9 \times 0.85 = 76.5\%$.  This
error compounding is the fundamental challenge of pipeline IE systems.

\subsection{Event Extraction}

Event extraction goes beyond entity-relation triples to extract structured event
representations:

\begin{quote}
\emph{``Apple acquired Beats Electronics for \$3 billion in May 2014.''}\\
Event type: \textsc{Acquisition}\\
Trigger: ``acquired''\\
Arguments: \{buyer: Apple, target: Beats Electronics, price: \$3B, date: May 2014\}
\end{quote}

The task decomposes into: (1) trigger identification (which word/phrase triggers the event),
(2) event type classification, and (3) argument extraction and role labeling.

\subsection{Coreference Resolution}

Coreference resolution groups mentions in a document that refer to the same real-world entity.
This is essential for IE because relations often span sentences:

\begin{quote}
\emph{``Tim Cook joined Apple in 1998.  He became CEO in 2011.''}
\end{quote}

Without resolving ``He'' $\to$ ``Tim Cook,'' the relation (\textit{Tim Cook},
\textsc{became\_ceo}, \textit{Apple}) cannot be extracted from the second sentence alone.

Modern coreference systems use mention-pair scoring: for each pair of spans, a model (typically
BERT-based) predicts whether they are coreferent.  The scores are used to build clusters via
agglomerative clustering or a span-ranking approach (Lee et al., 2017).

\subsection{End-to-End IE with LLMs}

LLMs can perform the entire IE pipeline in a single pass via structured prompting:

\begin{lstlisting}[language=Python, caption={LLM-based end-to-end information extraction.}]
prompt = """Extract all entities, relationships, and events
from the following text. Return as JSON.

Text: "Apple acquired Beats Electronics for $3B in May 2014.
Tim Cook announced the deal at WWDC."

Output format:
{"entities": [...], "relations": [...], "events": [...]}"""
\end{lstlisting}

The advantages are: no pipeline error compounding, handles implicit relations and events, and
adapts to new schemas via prompt modification.  The disadvantages are: inconsistent output
structure, higher latency, and difficulty in controlling precision/recall tradeoffs.  For
production systems, a hybrid approach---LLM for complex or ambiguous cases, dedicated models
for high-volume extraction---typically offers the best balance.

% ============================================================================
\section{Task Selection Decision Framework}
\label{sec:task_decision_framework}
% ============================================================================

The following table summarizes when to use each approach tier across the tasks covered in
this chapter.  The right choice depends on the constraint set, not on which approach is
``newest.''

\begin{table}[H]
\centering
\caption{Approach selection by task and constraint.  ``Classical'' = feature-based ML,
``Neural'' = fine-tuned BERT/RoBERTa, ``LLM'' = prompted large language model.}
\label{tab:task_approach_selection}
\small
\begin{tabular}{@{}lcccc@{}}
\toprule
\textbf{Constraint} & \textbf{Classical} & \textbf{Neural (FT)} & \textbf{LLM (Prompt)} \\
\midrule
Labeled data $<$ 100          & Good      & Poor     & Best     \\
Labeled data $>$ 1K           & Good      & Best     & Overkill \\
Latency $<$ 5 ms              & Best      & Fair     & Poor     \\
Schema changes frequently     & Poor      & Poor     & Best     \\
Strict output structure       & Best      & Good     & Fair     \\
Interpretability required     & Best      & Fair     & Poor     \\
Multilingual (50+ languages)  & Poor      & Good     & Best     \\
Cost per prediction matters   & Best      & Good     & Poor     \\
\bottomrule
\end{tabular}
\end{table}

\begin{warningbox}
The strongest interview signal on applied NLP questions is \textbf{not} knowing the most
advanced approach---it is knowing \emph{which} approach to use and \emph{why}.  Proposing an
LLM for a spam filter processing 10 million emails per day shows lack of engineering judgment.
Proposing Naive Bayes for a nuanced legal document analysis task shows lack of awareness of
model capabilities.  The right answer starts with constraints (data, latency, cost, accuracy
requirements) and derives the approach from them.
\end{warningbox}

% ============================================================================
\section{Interview Questions}
\label{sec:practical_tasks_interview}
% ============================================================================

\begin{interviewq}
\textbf{Q: How would you build a text classification system? Walk through the full pipeline
from data to deployment.}

\badgeLfive{L5} \quad \badgetype{System Design} \quad \badgefreq{\freqhigh}

\stronganswer
I would follow a structured pipeline with increasing complexity, stopping when the accuracy
meets requirements:

\textbf{1. Problem definition:} Clarify the label taxonomy (how many classes, single-label vs.
multi-label, hierarchical?), data availability, latency requirements, and accuracy targets.

\textbf{2. Data pipeline:} Collect and clean training data.  Perform label quality audit (what
is the inter-annotator agreement?).  Split into train/validation/test with stratification.
Ensure the test set is representative of production distribution.

\textbf{3. Baseline (day 1):} TF-IDF (unigrams + bigrams) + logistic regression.  This
trains in seconds, establishes a performance floor, and reveals whether the task is easy
(baseline $>$ 90\%) or hard (baseline $<$ 70\%).

\textbf{4. Stronger model:} Fine-tune BERT or DistilBERT on the training set.  Compare
against the baseline.  If the improvement is marginal, the baseline may be sufficient.

\textbf{5. Error analysis:} Examine misclassified examples.  Are errors due to ambiguous
labels, insufficient data for rare classes, domain mismatch, or model limitations?  Fix data
issues before model issues.

\textbf{6. Deployment:} Choose the model that meets the accuracy-latency-cost Pareto
frontier.  Implement monitoring: track prediction distribution shifts, confidence score
distributions, and flagged low-confidence predictions for human review.

\redflags
\begin{itemize}
    \item Jumping straight to BERT or LLMs without establishing a baseline.
    \item Not discussing data quality, label definitions, or evaluation methodology.
    \item Ignoring latency and cost constraints.
    \item Not mentioning error analysis or monitoring.
\end{itemize}

\interviewtest{End-to-end system design thinking, engineering maturity, understanding that
data quality often matters more than model choice.}

\goodgreat
\begin{itemize}
    \item \badgeLfive{L5}: Describes the baseline-then-upgrade approach, includes data
    preprocessing, and mentions evaluation with train/val/test split.
    \item \badgeLsix{L6}: Adds active learning for labeling, discusses the full Pareto
    frontier of model choices (NB $\to$ SVM $\to$ DistilBERT $\to$ BERT $\to$ LLM),
    includes deployment monitoring and A/B testing against rule-based fallbacks.
    \item \badgeLseven{L7}: Discusses how to handle evolving label taxonomies, multi-model
    architectures (fast model for easy cases, expensive model for hard cases), data flywheel
    (model predictions $\to$ human review $\to$ training data), and when to use LLM-based
    labeling to bootstrap the training set.
\end{itemize}

\followups{How would you handle class imbalance? What if the label taxonomy changes quarterly?
How would you decide between DistilBERT and full BERT for production? How would you monitor
for data drift?}
\end{interviewq}

\vspace{0.5em}

% --- Question 2 ---
\begin{interviewq}
\textbf{Q: How would you build a production NER system for a new domain (e.g., biomedical)
with limited labeled data?}

\badgeLfive{L5} \quad \badgetype{System Design} \quad \badgefreq{\freqhigh}

\stronganswer
This is a cold-start problem with a well-known playbook:

\textbf{Step 1: Define the entity schema.}  Work with domain experts to define entity types,
annotation guidelines, and ambiguous cases.  For biomedical NER: diseases, drugs, genes,
dosages, procedures.  Clear guidelines are essential---inconsistent annotations waste
everything downstream.

\textbf{Step 2: Bootstrap with existing resources.}  Use transfer learning from general NER
(fine-tune a model pretrained on CoNLL-2003) or, better, from a domain-pretrained model
(BioBERT, PubMedBERT for biomedical text).  Domain-specific pretraining captures vocabulary
and distributional patterns that general models miss.

\textbf{Step 3: Leverage gazetteers.}  Compile lists of known entities from domain knowledge
bases (UMLS for medical, SEC filings for financial).  Use these for: (a) weak labeling
(dictionary matching provides noisy labels for training), (b) feature augmentation (binary
feature indicating gazetteer membership), and (c) evaluation (check if the model finds
entities not in the gazetteer).

\textbf{Step 4: Active learning.}  Train an initial model on weakly labeled + any available
labeled data.  Score unlabeled sentences by model uncertainty (entropy of per-token
predictions).  Prioritize the most uncertain sentences for expert annotation.  This
maximizes the value of each annotation dollar.

\textbf{Step 5: Iterate.}  Retrain with new annotations, evaluate on a held-out test set,
perform error analysis, and repeat.  50--200 carefully annotated documents often suffice to
reach production quality when combined with transfer learning and gazetteers.

\redflags
\begin{itemize}
    \item Proposing to ``just train BERT on a few examples'' without transfer learning or
    data augmentation strategy.
    \item Not mentioning domain-specific pretrained models.
    \item Ignoring the importance of annotation guidelines and inter-annotator agreement.
    \item Not discussing gazetteers or weak supervision.
\end{itemize}

\interviewtest{Ability to design a practical system under data constraints, knowledge of
transfer learning and active learning, domain adaptation awareness.}

\goodgreat
\begin{itemize}
    \item \badgeLfive{L5}: Describes transfer learning from domain-pretrained models, mentions
    gazetteers, and outlines an iterative annotation-training loop.
    \item \badgeLsix{L6}: Adds weak supervision frameworks (Snorkel-style labeling functions),
    discusses the CRF layer decision, and mentions evaluation at the entity level (exact match
    F1 vs.\ partial match vs.\ type-specific F1).
    \item \badgeLseven{L7}: Discusses continued pretraining on domain corpus before
    fine-tuning (two-stage adaptation), few-shot NER with LLMs for initial labeling, and
    how to handle nested entities and entity linking in the domain.
\end{itemize}

\followups{How do you handle nested entities? What is the difference between exact and partial
match F1 for NER evaluation? How would you use an LLM to accelerate annotation?}
\end{interviewq}

\vspace{0.5em}

% --- Question 3 ---
\begin{interviewq}
\textbf{Q: Compare extractive vs.\ abstractive summarization. What are the tradeoffs and when
would you use each?}

\badgeLfive{L5} \quad \badgetype{Trade-off} \quad \badgefreq{\freqhigh}

\stronganswer
The core tradeoff is \textbf{faithfulness vs.\ fluency}.

\textbf{Extractive summarization} selects sentences verbatim from the source.  Advantages:
zero hallucination risk (every sentence is a direct quote), computationally cheaper,
interpretable (can trace every sentence to its source).  Disadvantages: summaries are often
choppy and redundant, cannot paraphrase or compress, and include irrelevant details from
selected sentences.

\textbf{Abstractive summarization} generates new text that may paraphrase or fuse information.
Advantages: more natural and concise, can synthesize across multiple passages, produces
coherent narratives.  Disadvantages: prone to hallucination (both intrinsic and extrinsic),
harder to verify, and requires more compute.

\textbf{Decision framework:}
\begin{itemize}
    \item High-stakes domains (medical, legal, financial): prefer extractive or
    extractive-then-polish (extract key sentences, then use an LLM to smooth the language
    while grounding each claim in a source sentence).
    \item User-facing summaries where readability matters: abstractive with faithfulness
    checks (NLI-based verification, citation generation).
    \item Long documents: extractive first (sentence ranking to identify key content), then
    abstractive on the selected sentences.
\end{itemize}

\redflags
\begin{itemize}
    \item Saying abstractive is ``always better because it generates more natural text.''
    \item Not mentioning hallucination as the central risk of abstractive summarization.
    \item Not discussing faithfulness evaluation.
    \item Treating this as a binary choice rather than a spectrum.
\end{itemize}

\interviewtest{Understanding of practical tradeoffs in generation tasks, awareness of
hallucination risk, ability to match approach to constraints.}

\goodgreat
\begin{itemize}
    \item \badgeLfive{L5}: Clearly states the faithfulness-fluency tradeoff and gives
    concrete scenarios for each approach.
    \item \badgeLsix{L6}: Discusses hybrid approaches (extract-then-abstract), faithfulness
    evaluation methods (NLI, QuestEval, FactCC), and ROUGE's blindness to factual accuracy.
    \item \badgeLseven{L7}: Discusses the production pipeline with guardrails (faithfulness
    classifier as a post-generation filter), how to handle long documents that exceed context
    windows, and the role of retrieval-augmented generation for grounded summarization.
\end{itemize}

\followups{How would you detect hallucinations in a generated summary? What is ROUGE measuring
and what does it miss? How would you summarize a 100-page legal document?}
\end{interviewq}

\vspace{0.5em}

% --- Question 4 ---
\begin{interviewq}
\textbf{Q: Design a question answering system for a company's internal documentation
(Confluence, Notion, internal wikis).}

\badgeLsix{L6} \quad \badgetype{System Design} \quad \badgefreq{\freqmed}

\stronganswer
This is a RAG (Retrieval-Augmented Generation) system design problem with specific constraints:
the corpus is private (cannot use external search), documents update frequently, and accuracy
is critical (employees act on the answers).

\textbf{Ingestion pipeline:}
\begin{enumerate}
    \item \textbf{Document processing:} Parse documents from various formats (Confluence
    Markdown, Notion blocks, PDF, HTML).  Extract text, tables, code blocks, and metadata
    (author, last updated, access permissions).
    \item \textbf{Chunking:} Split documents into retrievable units.  Use semantic chunking
    (split at section boundaries, paragraphs) rather than fixed-size chunks.  Overlap of
    50--100 tokens between chunks preserves context at boundaries.  Typical chunk size:
    256--512 tokens.
    \item \textbf{Embedding:} Encode chunks using a bi-encoder (e.g., E5, BGE, or a
    domain-fine-tuned model).  Store embeddings in a vector database (Pinecone, Weaviate,
    Qdrant).
    \item \textbf{Metadata indexing:} Also build a BM25 index over chunk text for hybrid
    retrieval.  Index metadata (document title, author, date) for filtering.
\end{enumerate}

\textbf{Query pipeline:}
\begin{enumerate}
    \item \textbf{Retrieval:} Hybrid retrieval (dense + BM25) with reciprocal rank fusion.
    Retrieve top-20 candidates.
    \item \textbf{Re-ranking:} Cross-encoder re-ranker scores each (query, chunk) pair.
    Select top-5.
    \item \textbf{Generation:} Feed the query and top-5 chunks to an LLM with a prompt that
    instructs it to answer based only on the provided context and cite sources.
    \item \textbf{Citation:} The response includes inline citations linking each claim to the
    source document.  This enables verification and builds trust.
\end{enumerate}

\textbf{Key design decisions:} Access control (only retrieve documents the user has permission
to see), freshness (re-index on document updates), handling ``I don't know'' (when the
retrieved context doesn't contain the answer, the system should say so), and feedback loop
(users can upvote/downvote answers to improve retrieval and generation over time).

\redflags
\begin{itemize}
    \item Designing a system that ignores document access permissions.
    \item Not including a ``no answer'' / ``I don't know'' pathway.
    \item Using only dense retrieval without BM25 (misses exact keyword matches).
    \item Not discussing chunking strategy or its impact on retrieval quality.
\end{itemize}

\interviewtest{End-to-end RAG system design, practical awareness of enterprise constraints
(permissions, freshness, trust), and understanding of the full retrieval-generation pipeline.}

\goodgreat
\begin{itemize}
    \item \badgeLfive{L5}: Describes the basic RAG pipeline with retrieval and generation,
    mentions chunking and embedding.
    \item \badgeLsix{L6}: Adds hybrid retrieval, re-ranking, citation, access control,
    incremental indexing, and evaluation strategy (retrieval recall, answer quality, user
    satisfaction).
    \item \badgeLseven{L7}: Discusses query routing (simple factoid queries to lightweight
    extraction, complex queries to full RAG), fine-tuning the embedding model on internal
    query logs, handling tables and structured data in documents, and a feedback-driven
    continuous improvement loop.
\end{itemize}

\followups{How would you handle tables in the documentation? What chunk size would you start
with and how would you tune it? How would you evaluate this system before and after launch?}
\end{interviewq}

\vspace{0.5em}

% --- Question 5 ---
\begin{interviewq}
\textbf{Q: How would you handle aspect-based sentiment analysis at scale?}

\badgeLfive{L5} \quad \badgetype{Conceptual} \quad \badgefreq{\freqmed}

\stronganswer
Aspect-based sentiment analysis (ABSA) extracts structured (aspect, opinion, sentiment)
triples from text.  At scale, the key challenge is balancing extraction quality with
throughput.

\textbf{Approach selection:}
For high-throughput production (millions of reviews per day), I would use a fine-tuned BERT
model that jointly performs aspect extraction and sentiment classification.  The model takes
the review as input and either: (a) performs sequence labeling to extract aspect and opinion
spans with BIO tags, then classifies sentiment per extracted pair, or (b) frames the task
generatively---fine-tune a T5 model to output structured triples directly (e.g.,
``(food, excellent, positive); (service, slow, negative)'').

\textbf{For cold-start or prototyping:} Use an LLM with a structured output prompt.  This
handles implicit aspects (``The wait was worth it'' implies positive sentiment about food
quality) better than extraction-based methods, but costs 100--1000x more per prediction.

\textbf{Handling implicit aspects:}  A review like ``overpriced for what you get'' expresses
negative sentiment about value, but ``value'' is never explicitly mentioned.  This requires
aspect category detection (mapping to predefined categories) rather than just span extraction.

\textbf{Domain adaptation:}  ABSA models trained on restaurant reviews fail on electronics
reviews because the aspect vocabulary is entirely different.  Solutions: domain-specific
fine-tuning, aspect category ontology per domain, and shared sentiment layers with
domain-specific aspect layers.

\redflags
\begin{itemize}
    \item Treating ABSA as simple sentence-level sentiment classification.
    \item Not distinguishing explicit from implicit aspects.
    \item Proposing only an LLM solution without discussing cost at scale.
\end{itemize}

\interviewtest{Understanding of structured prediction tasks, ability to decompose a complex
NLP task into subtasks, and awareness of the explicit vs.\ implicit aspect distinction.}

\goodgreat
\begin{itemize}
    \item \badgeLfive{L5}: Correctly decomposes ABSA into subtasks, mentions both
    extraction-based and generative approaches.
    \item \badgeLsix{L6}: Discusses implicit aspects, domain adaptation, and the
    throughput-quality tradeoff between fine-tuned models and LLMs.
    \item \badgeLseven{L7}: Discusses multi-task learning (joint aspect-sentiment models),
    cross-lingual ABSA for global products, aggregation of aspect-level sentiment into
    product-level insights, and using ABSA output to drive business actions (e.g., routing
    complaints by aspect category to the relevant team).
\end{itemize}

\followups{How would you aggregate aspect-level sentiment across thousands of reviews into
actionable product insights? How would you handle aspects that are mentioned with abbreviations
or slang?}
\end{interviewq}

\vspace{0.5em}

% --- Question 6 ---
\begin{interviewq}
\textbf{Q: Compare Naive Bayes vs.\ fine-tuned BERT for text classification. When might Naive
Bayes win?}

\badgeLfive{L5} \quad \badgetype{Trade-off} \quad \badgefreq{\freqmed}

\stronganswer
Naive Bayes wins in specific, well-defined scenarios:

\textbf{1. Extremely small training data ($<$ 100 examples):}  Naive Bayes has
$O(|V| \times |C|)$ parameters---far fewer than BERT's 110M.  With 50 examples, BERT's
classification head overfits, while Naive Bayes's strong inductive bias (conditional
independence) acts as regularization.  The independence assumption is a ``useful wrong
assumption'' that prevents overfitting in small-data regimes.

\textbf{2. Hard latency constraints ($<$ 1 ms):}  Naive Bayes inference is a sparse dot
product---microseconds.  BERT requires a full forward pass through 12 transformer layers---
10--50 ms on GPU, hundreds of ms on CPU.  For high-QPS systems (email spam filtering at
100K+ requests per second), this difference is decisive.

\textbf{3. Interpretability:}  Naive Bayes feature log-ratios directly explain predictions.
``This email was flagged as spam because `free,' `winner,' and `click' have high spam
likelihood ratios.''  BERT is opaque without additional interpretability infrastructure.

\textbf{4. No GPU infrastructure:}  Naive Bayes trains in seconds on a laptop CPU and
requires no GPU for inference.  BERT fine-tuning and inference are impractical without GPU
access.

\textbf{When BERT wins:}  When data is sufficient ($>$ 1K examples), when the task requires
understanding context, negation, or compositional semantics, and when latency budgets
permit.  For most classification tasks with adequate data, BERT outperforms Naive Bayes by
5--15\% absolute accuracy.

\redflags
\begin{itemize}
    \item Saying ``always use BERT'' or ``Naive Bayes is obsolete.''
    \item Not mentioning latency or resource constraints.
    \item Not quantifying the data threshold where BERT overtakes Naive Bayes.
\end{itemize}

\interviewtest{Engineering judgment, practical awareness of the simplicity-performance
tradeoff, ability to reason about constraints rather than defaulting to the most powerful
model.}

\goodgreat
\begin{itemize}
    \item \badgeLfive{L5}: Gives at least two concrete scenarios where Naive Bayes wins, with
    clear reasoning.
    \item \badgeLsix{L6}: Discusses the full model hierarchy
    (NB $\to$ SVM $\to$ DistilBERT $\to$ BERT $\to$ LLM), mentions Pareto frontier of
    accuracy vs.\ latency, and proposes DistilBERT as a middle ground.
    \item \badgeLseven{L7}: Discusses cascade architectures (fast model for easy cases,
    expensive model for hard cases), cost-benefit analysis for choosing models, and how
    LLM-generated labels can bootstrap training data for Naive Bayes when no labeled data
    exists.
\end{itemize}

\followups{What about DistilBERT as a middle ground? How would you quantify the business
value of the accuracy gain from BERT? What about few-shot LLM classification as an
alternative when no training data exists?}
\end{interviewq}

\vspace{0.5em}

% --- Question 7 ---
\begin{interviewq}
\textbf{Q: How would you detect sarcasm in text? Why is this hard?}

\badgeLsix{L6} \quad \badgetype{Conceptual} \quad \badgefreq{\freqlow}

\stronganswer
Sarcasm detection is hard because it requires understanding what the speaker \emph{means}
versus what they \emph{say}---a form of pragmatic inference that goes beyond surface
semantics.

\textbf{Why it is fundamentally hard:}
\begin{itemize}
    \item \textbf{Surface-meaning inversion:} ``Great, another Monday'' uses positive words
    to express negative sentiment.  Models that rely on lexical sentiment are systematically
    wrong on sarcasm.
    \item \textbf{Context dependence:} ``Nice weather'' is literal on a sunny day, sarcastic
    during a hurricane.  The text alone is insufficient---the model needs situational context.
    \item \textbf{Cultural and personal variation:} Sarcasm conventions vary across cultures,
    age groups, and individuals.  What is sarcastic in British English may be literal in
    another context.
    \item \textbf{Absence of prosodic cues:} In speech, sarcasm is often conveyed through
    tone, emphasis, and pacing.  Text strips these signals away.
\end{itemize}

\textbf{Approaches:}
\begin{itemize}
    \item \textbf{Feature-based:} Incongruity features (positive words + negative context),
    punctuation patterns (excessive exclamation marks, ellipsis), capitalization (``SURE''),
    and discourse features (contrast between sentences).
    \item \textbf{Contextual models:} Fine-tuned BERT on sarcasm-annotated datasets (Twitter
    sarcasm corpus, Reddit \texttt{/s} tag data).  Achieves 70--80\% F1, well below human
    performance.
    \item \textbf{Multimodal:} Incorporating images (sarcastic memes), audio (tone of voice),
    or video.  Multimodal approaches show significant improvements but are limited to
    settings where these modalities are available.
    \item \textbf{LLMs:} Better at sarcasm than previous models due to broader world knowledge,
    but still unreliable, especially on subtle or culture-specific sarcasm.
\end{itemize}

\textbf{Practical implication:} In production sentiment analysis, explicitly model an
``uncertain/mixed'' category rather than forcing every sarcastic text into positive or
negative.  Flag potentially sarcastic texts (high positive-word count + negative context
signals) for human review.

\redflags
\begin{itemize}
    \item Claiming that a large enough model ``solves'' sarcasm.
    \item Not acknowledging the fundamental role of non-textual context.
    \item Proposing a purely lexical approach (keyword lists for sarcasm detection).
\end{itemize}

\interviewtest{Understanding of pragmatic language phenomena, intellectual honesty about
the limits of current NLP systems, and ability to design systems that gracefully handle
unsolvable subproblems.}

\goodgreat
\begin{itemize}
    \item \badgeLfive{L5}: Explains why sarcasm is hard (surface inversion, context
    dependence) and names at least one approach.
    \item \badgeLsix{L6}: Discusses the role of multimodal cues, the limits of text-only
    detection, and practical strategies for handling sarcasm in production pipelines
    (uncertainty, flagging, human-in-the-loop).
    \item \badgeLseven{L7}: Connects to pragmatics research (Grice's maxims, conversational
    implicature), discusses how sarcasm detection could benefit from user modeling
    (individual sarcasm tendency), and acknowledges the natural ceiling on text-only sarcasm
    detection accuracy.
\end{itemize}

\followups{How would you collect sarcasm training data? What is the inter-annotator agreement
for sarcasm labeling? How does sarcasm detection differ across languages?}
\end{interviewq}

\vspace{0.5em}

% --- Question 8 ---
\begin{interviewq}
\textbf{Q: Design a task-oriented dialogue system for restaurant booking.}

\badgeLsix{L6} \quad \badgetype{System Design} \quad \badgefreq{\freqmed}

\stronganswer
I would design this as a modular system with clear interfaces between components, allowing
independent iteration on each.

\textbf{1. NLU Module:}
\begin{itemize}
    \item \textbf{Intent classification:} Intents include \texttt{book\_table},
    \texttt{cancel\_reservation}, \texttt{modify\_reservation}, \texttt{check\_hours},
    \texttt{ask\_menu}, \texttt{general\_inquiry}.  Fine-tuned BERT or DistilBERT on
    intent-labeled utterances.
    \item \textbf{Slot filling:} Slots include \texttt{cuisine}, \texttt{location},
    \texttt{date}, \texttt{time}, \texttt{party\_size}, \texttt{restaurant\_name},
    \texttt{price\_range}.  Joint intent-slot model using BERT with a sequence labeling
    head for slots and a classification head for intent on \texttt{[CLS]}.
\end{itemize}

\textbf{2. Dialogue State Tracker:}
Maintains a belief state (slot-value pairs) accumulated across turns.  Must handle slot
updates (``Actually, make that 6 people''), slot confirmation (``Yes, 7pm is correct''), and
slot negation (``Not Italian, I want Thai'').  Implementation: rule-based tracker on top of
NLU output for reliability, or a BERT-based generative DST model for flexibility.

\textbf{3. Dialogue Policy:}
Given the belief state and available API results, determine the next system action:
\begin{itemize}
    \item If required slots are missing: ask for the missing slot.
    \item If all required slots are filled: query the restaurant database.
    \item If results are available: present options and ask for confirmation.
    \item If confirmed: make the reservation and confirm to the user.
\end{itemize}
Start with a rule-based policy (deterministic and predictable), evolve to a learned policy if
the conversation flows become too complex for rules.

\textbf{4. NLG Module:}
Template-based for reliability: ``I found \{\texttt{n}\} \{\texttt{cuisine}\} restaurants near
\{\texttt{location}\}.  Would you like to book at \{\texttt{restaurant\_name}\} for
\{\texttt{party\_size}\} at \{\texttt{time}\}?''  Optionally backed by an LLM for more
natural phrasing, with template fallback if the LLM output fails validation.

\textbf{5. Error recovery:}
Handle NLU failures gracefully: ``I didn't catch the time.  When would you like to dine?''
Implement maximum retry limits and escalation to human agent.

\redflags
\begin{itemize}
    \item Proposing ``just use ChatGPT'' without discussing controllability, reliability,
    or the structured nature of the task.
    \item Not defining the slot schema or intent taxonomy.
    \item Ignoring error recovery and edge cases.
    \item Not discussing how the system interacts with the restaurant database/API.
\end{itemize}

\interviewtest{Ability to design a structured, multi-component system with clear interfaces,
understanding of task-oriented dialogue architecture, and practical awareness of reliability
requirements.}

\goodgreat
\begin{itemize}
    \item \badgeLfive{L5}: Describes the NLU $\to$ DST $\to$ Policy $\to$ NLG pipeline with
    correct terminology and a concrete slot schema.
    \item \badgeLsix{L6}: Discusses hybrid approaches (rule-based policy for reliability,
    learned NLU for flexibility), error recovery, evaluation metrics (task completion rate,
    average turns to completion, slot accuracy), and the database interaction.
    \item \badgeLseven{L7}: Discusses when to transition from a pipeline to an end-to-end LLM
    approach (and the controllability tradeoffs), multi-language support, personalization
    (remembering user preferences across sessions), and A/B testing the dialogue policy.
\end{itemize}

\followups{How would you evaluate this system? What metrics matter? How would you handle
ambiguous user inputs? When would you switch from a pipeline to an end-to-end LLM-based
system?}
\end{interviewq}

\vspace{0.5em}

% --- Question 9 ---
\begin{interviewq}
\textbf{Q: You need to build an information extraction system that populates a knowledge graph
from news articles. Walk through your approach.}

\badgeLsix{L6} \quad \badgetype{System Design} \quad \badgefreq{\freqmed}

\stronganswer
I would build a pipeline IE system with four stages, with careful attention to error
compounding across stages.

\textbf{Stage 1: Entity Extraction.}
Fine-tune a BERT-CRF model for NER on news text (CoNLL-2003 or a custom-annotated dataset).
Entity types: Person, Organization, Location, Date, Event, Product.  Use BIOES tagging for
better boundary detection.  Expected F1: 88--93\% depending on domain.

\textbf{Stage 2: Coreference Resolution.}
Cluster entity mentions that refer to the same entity.  Use a span-based coreference model
(e.g., based on Lee et al., 2017, fine-tuned on OntoNotes).  This is critical: without
coreference, the system misses relations expressed across sentences (``Tim Cook...He became
CEO...'').

\textbf{Stage 3: Relation Extraction.}
For each pair of entities within a document (or within a window), classify the relation.
Use BERT with typed entity markers.  Define a relation schema aligned with the target KG
ontology (e.g., Wikidata properties: CEO\_of, headquartered\_in, founded\_in, acquired).
Include a ``no relation'' class---most entity pairs have no meaningful relation.

\textbf{Stage 4: Entity Linking + KG Population.}
Link extracted entities to canonical KB entries using candidate generation (string matching
+ dense retrieval) and candidate ranking (cross-encoder).  Insert validated triples into
the knowledge graph.

\textbf{Error mitigation:}  Since errors compound across stages ($0.9^4 = 65.6\%$ for four
stages at 90\% each), I would: (a) use confidence thresholds at each stage and only pass
high-confidence outputs forward, (b) implement a human-in-the-loop verification for
low-confidence triples, and (c) consider an end-to-end LLM approach for complex or ambiguous
cases that the pipeline handles poorly.

\redflags
\begin{itemize}
    \item Not mentioning coreference resolution (a critical and often-forgotten stage).
    \item Ignoring error compounding across pipeline stages.
    \item Not defining a relation schema or entity types.
    \item Proposing only an LLM solution without discussing throughput or cost for large
    document volumes.
\end{itemize}

\interviewtest{Understanding of the full IE pipeline, awareness of error compounding,
practical system design under real-world constraints.}

\goodgreat
\begin{itemize}
    \item \badgeLfive{L5}: Describes NER $\to$ RE $\to$ entity linking as the core pipeline
    and mentions evaluation per stage.
    \item \badgeLsix{L6}: Adds coreference resolution, discusses error compounding with
    quantitative analysis, and proposes confidence-based filtering and human-in-the-loop
    verification.
    \item \badgeLseven{L7}: Discusses incremental KG updates (handling entity merges, relation
    updates, temporal validity), combines pipeline with LLM-based extraction for long-tail
    relations, and designs a feedback loop where KG consumers report errors that improve
    extraction models.
\end{itemize}

\followups{How would you handle contradictory information from different sources? How would
you evaluate the quality of the populated knowledge graph? How would you handle temporal
relations (``was CEO from 2011 to 2023'')?}
\end{interviewq}

\vspace{0.5em}

% --- Question 10 ---
\begin{interviewq}
\textbf{Q: When is an LLM overkill for an NLP task? Give concrete examples.}

\badgeLfive{L5} \quad \badgetype{Conceptual} \quad \badgefreq{\freqhigh}

\stronganswer
An LLM is overkill whenever a simpler approach achieves acceptable accuracy within the system's
constraints.  Concrete examples:

\textbf{1. Language detection:} Identifying whether text is English, French, or Chinese
requires no semantic understanding.  A character n-gram model (fastText's language ID, or
even simple Unicode script detection) achieves $>$99\% accuracy in microseconds.  An LLM
would be 10,000x slower and no more accurate.

\textbf{2. High-volume spam filtering:} Email spam filters process billions of messages daily.
TF-IDF + logistic regression achieves 95--98\% accuracy with sub-millisecond latency and
trivial infrastructure cost.  An LLM at \$0.01 per 1K tokens would cost millions of dollars
per day with no meaningful accuracy improvement on this well-defined binary task.

\textbf{3. Keyword-based routing:} Routing customer support tickets by category (billing,
technical, shipping) based on keyword presence.  A simple rule-based system or TF-IDF
classifier handles this perfectly; the categories are well-defined and the signal is in
obvious keywords.

\textbf{4. Regex-solvable extraction:} Extracting phone numbers, email addresses, dates, or
URLs from text.  Regular expressions are exact, fast, and have zero false positives when
correctly written.  An LLM would be slower, more expensive, and less reliable.

\textbf{5. Deduplication:} Near-duplicate document detection using MinHash or SimHash.  This
is a similarity computation problem, not a language understanding problem.

\textbf{The general principle:} Use the simplest model that meets accuracy requirements within
the latency and cost budget.  LLMs are appropriate when the task requires genuine language
understanding, reasoning, or generation---and when simpler approaches have been tried and found
insufficient.

\redflags
\begin{itemize}
    \item Unable to name concrete examples---gives only vague statements about ``simple tasks.''
    \item Saying ``always start with an LLM and optimize later.''
    \item Not considering cost or latency in the analysis.
\end{itemize}

\interviewtest{Engineering judgment, cost awareness, ability to resist hype and select
appropriate tools for the problem.}

\goodgreat
\begin{itemize}
    \item \badgeLfive{L5}: Names 3+ concrete examples with clear reasoning about why simpler
    approaches suffice.
    \item \badgeLsix{L6}: Quantifies the cost and latency differences, discusses the model
    selection hierarchy (rules $\to$ classical ML $\to$ fine-tuned BERT $\to$ LLM), and
    frames this as a Pareto optimization problem.
    \item \badgeLseven{L7}: Discusses when LLMs are the right starting point
    (rapid prototyping, undefined label spaces, tasks requiring world knowledge or reasoning),
    and how to decide when to ``graduate'' from LLM to dedicated model (when volume and cost
    justify the engineering investment in training and deploying a task-specific model).
\end{itemize}

\followups{When IS an LLM the right choice? How would you decide when to switch from an LLM
prototype to a dedicated model? How would you estimate the cost of an LLM-based solution vs.\
a classical one?}
\end{interviewq}

\bigskip

\begin{keyinsight}[Chapter Summary: What Interviewers Are Really Testing]
When interviewers ask about practical NLP tasks, they are testing three things beyond technical
knowledge: (1) \emph{Can you decompose a complex problem into well-defined subtasks?}  The
best candidates break down ``build a dialogue system'' into NLU, DST, policy, and NLG with
clear interfaces.  (2) \emph{Do you have engineering judgment about approach selection?}  The
answer is never ``use the best model''---it is ``use the right model given the constraints.''
(3) \emph{Do you think about the full system, not just the model?}  Data collection, annotation
guidelines, evaluation metrics, monitoring, error analysis, and deployment are as important as
model architecture.  Candidates who discuss all of these dimensions signal that they have
actually built and shipped NLP systems, not just trained models on benchmarks.
\end{keyinsight}
